%! Parametrization of Implicitly Defined Functions — Unit 4, Lesson 4.7 — Homework
\documentclass[10pt]{article}
\usepackage{precalculus-article}
\usepackage{precalculus-boxes}
\newcommand{\TallMath}[1]{$\displaystyle #1\rule[-1.4em]{0pt}{3.2em}$}

\begin{document}

\pageheader{Unit 4, Lesson 4.7}{Homework}
\namedateperiod

\vspace{0.10in}

\begin{notesbox}{Parametrization Practice}

\textbf{1.}\quad The ellipse $\dfrac{x^2}{9} + \dfrac{y^2}{25} = 1$ is centered at the
origin with $a = 3$ and $b = 5$.

\vspace{4pt}
(a)\quad Write the parametric equations that trace this ellipse counterclockwise, $0 \le t \le 2\pi$.

\vspace{4pt}
$x(t) = $ \blank{4cm} \qquad $y(t) = $ \blank{4cm}

\vspace{10pt}
(b)\quad Find the coordinates of the four extreme points of the ellipse using your parametrization.
Show the $t$-value for each.

\vspace{4pt}
\begin{center}
\renewcommand{\arraystretch}{1.6}
\begin{tabular}{|l|c|c|}
\hline
\textbf{Extreme point} & $t$ & \textbf{Coordinates} \\
\hline
Rightmost & \blank{2cm} & \blank{3cm} \\
\hline
Leftmost  & \blank{2cm} & \blank{3cm} \\
\hline
Topmost   & \blank{2cm} & \blank{3cm} \\
\hline
Bottommost & \blank{2cm} & \blank{3cm} \\
\hline
\end{tabular}
\end{center}

\vspace{10pt}
\textbf{2.}\quad The ellipse $\dfrac{(x-3)^2}{16} + \dfrac{(y+1)^2}{4} = 1$.

\vspace{4pt}
(a)\quad Identify $h$, $k$, $a$, $b$, then write the parametric equations.
\writelines{2}

\vspace{4pt}
(b)\quad Use your parametrization to find all $t$-values where the ellipse crosses the
$x$-axis (set $y(t) = 0$ and solve). Then state the coordinates of those points.
\writelines{3}

\end{notesbox}

\vspace{0.08in}

\begin{notesbox}{Parabola Parametrization}

\textbf{3.}\quad The parabola $y = 2(x - 1)^2 - 3$. Let $x(t) = 1 + t$, so $y(t) = 2t^2 - 3$.

\vspace{4pt}
(a)\quad Find the $x$-intercepts of the parabola: set $y(t) = 0$ and solve for $t$.
Then find the corresponding $x$-values.
\writelines{3}

\vspace{4pt}
(b)\quad What is the vertex of the parabola? Use $t = 0$ in the parametrization to confirm.
\writelines{2}

\vspace{10pt}
\textbf{4.}\quad The circle $x^2 + y^2 = 49$.

\vspace{4pt}
(a)\quad Write the parametric equations. \quad $x(t) = $ \blank{3cm} \quad $y(t) = $ \blank{3cm}

\vspace{8pt}
(b)\quad Find the coordinates when $t = \pi/3$. (Use $\cos(\pi/3) = 1/2$, $\sin(\pi/3) = \sqrt{3}/2$.)
\writelines{2}

\end{notesbox}

\vspace{0.08in}

\begin{notesbox}{AP-Style Practice}

\textbf{5.}\quad Which of the following parametrizations correctly traces the ellipse
$\dfrac{x^2}{25} + \dfrac{y^2}{9} = 1$ counterclockwise for $0 \le t \le 2\pi$?

\vspace{6pt}
\begin{enumerate}[label=(\Alph*), leftmargin=2em, itemsep=5pt]
  \item $x(t) = 9\cos t,\quad y(t) = 25\sin t$
  \item $x(t) = 5\cos t,\quad y(t) = 3\sin t$
  \item $x(t) = 25\cos t,\quad y(t) = 9\sin t$
  \item $x(t) = 3\cos t,\quad y(t) = 5\sin t$
\end{enumerate}

\end{notesbox}

\end{document}
